\documentclass[10pt,a4paper]{amsart}
\usepackage[margin=19mm]{geometry}
\usepackage{amsmath,amssymb,mathtools}
\usepackage[scr=rsfs,cal=euler]{mathalfa}
\usepackage{xfrac}
\usepackage[unicode,hidelinks,bookmarks=false]{hyperref}
\newcommand{\ie}{i.e.\ }
\newcommand{\cf}{cf.\ }
\newcommand{\resp}{respectively\ }
\newcommand{\Subsection}{\subsection}
\providecommand{\nobreakdash}{\nobreak\hbox{-}}
\setlength{\emergencystretch}{2em}
\multlinegap0pt
\usepackage{amssymb}
\usepackage{bm}
\newtheorem{proposition}{Proposition}
\renewcommand{\H}{\mathcal H}
\newcommand{\N}{\mathbb{N}}
\newcommand{\T}{\mathbb{T}}
\newcommand{\dl}{\delta}
\renewcommand{\l}{\ell}
\newcommand{\les}{\lesssim}
\title{Nowhere continuity of the flow map of an integrable derivative nonlinear Schr\"odinger system on the torus}
\author{Toshiki Kondo}
\date{Source: arXiv:2607.05796v1 (CC BY 4.0)}
\begin{document}
\maketitle

\noindent\emph{Source and licence.} Nowhere continuity of the flow map of an integrable derivative nonlinear Schr\"odinger system on the torus. Version arXiv:2607.05796v1 (CC BY 4.0), distributed by arXiv under the Creative Commons Attribution 4.0 International licence, http://creativecommons.org/licenses/by/4.0/, as recorded in the arXiv licence metadata for this exact version and shown on the abstract page https://arxiv.org/abs/2607.05796v1. Full licence notice: \texttt{licenses/arxiv-2607.05796-CC-BY-4.0.txt}.\par\smallskip

\noindent\textit{Editorial scope.} Counted unit: the article's proposition prop:app (source0.tex lines 1226--1236). The article's complete original proof of it is delivered with it (source0.tex lines 1238--1440, one \texttt{proof} environment of 203 lines). No mathematics is written, repaired or re-derived: the statement and the proof are the article's own lines, in the article's own notation, and no retained line is substituted or re-flowed.  No further source-local context is needed: every cross-reference of every retained line resolves inside the two delivered spans.  Nothing was omitted from any retained span, so the package carries no omission record.  No retained line contains a citation, so the wrapper carries no bibliography and no bibliography entry.  No retained line contains a figure environment, an image-inclusion command or drawing code, so no image is delivered and the package declares no asset dependency.  Editorial formatting only, in addition: an AMS \texttt{amsart} preamble that restates the article's own declarations and macros used by the retained lines, namely the environments \texttt{proposition} on separate counters, together with the standard packages those lines need.  The retained environments print in this document as Proposition 1 in order of appearance, whereas the article prints the same items with the numbering of its own full document.  The complete native source of the article as distributed in the arXiv e-print for this exact version is bundled unchanged as \texttt{sources/204572/source0.tex}.
\begin{proposition}
\label{prop:app} 
Let $s > \frac 32$. For $0<\dl<1$ and $(u_0, v_0) \in \H^s(\T)$, there exists a sequence
$\{(u_{0,n}, v_{0,n})\}_{n \in \N} \subset \H^s(\T)$ such that the followings hold:
\begin{enumerate}
%\item $(u_{0,n}, v_{0,n}) \in \H^s(\T) \setminus \H^{s+1}(\T)$
\item $(u_{0,n}, v_{0,n})$ converges to $(u_0, v_0)$ in $\H^s(\T)$ as $n \to \infty$;
\item $\Im M (u_{0, n}, v_{0, n}) \neq 0$ for any $n \in \N$;
\item $(P_\pm u_{0,n}, P_\pm v_{0,n}) \notin \H^{s+\dl}(\T)$ for any $n \in \N$.
\end{enumerate}
\end{proposition}

\begin{proof}
We consider the two cases:
\begin{itemize}
\item Case 1: $\Im M(u_0, v_0) \neq 0$;
\item Case 2: $\Im M(u_0, v_0) = 0$.
\end{itemize}

Case 1: First, we consider the case $(u_0, v_0) \in \H^{s+\dl}(\T)$.
%\begin{itemize}
%\item Case 1-1: $(u_0, v_0) \in \H^{s+\dl}(\T)$
%\item Case 1-2: $(u_0, v_0) \notin \H^{s+\dl}(\T)$ 
%\end{itemize} 
Set
\[
u_{0}^{(\l)}(x) = u_0 (x) + \frac 1\l \sum_{k \in 2^\N} k^{-s-\dl} (e^{i k x} + e^{-ikx})
\]
for $\l \in \N$.
It follows from $u_0 \in H^{s+\dl}(\T)$ that
\begin{equation*}
P_\pm u_0^{(\l)} = P_\pm (u_0^{(\l)} - u_0) + P_\pm u_0 \notin H^{s+\dl}(\T)
%\label{app1}
\end{equation*}
for any $\l \in \N$.

A simple calculation yields that
\begin{equation}
\| u_{0}^{(\l)} - u_0 \|_{H^s}
= \bigg( \frac 2{\l^2}  \sum_{k \in 2^{\N}} k^{-2\dl} \bigg)^{\frac 12}
\les \frac 1\l
\label{app2}
\end{equation}
for any $\l \in \N$.
Note that
\begin{align*}
\Big|
\Im \int_\T u_0^{(\l)} v_0 dx
- \Im \int_\T u_0 v_0 dx
\Big|
\le
\| u_{0}^{(\l)} - u_0 \|_{L^2} \|v_0\|_{L^2}.
\end{align*}
From \eqref{app2} and $s>\frac 32$,
we obtain that
\[
\lim_{\l \to \infty} \Im \int_\T u_{0}^{(\l)} v_0 dx
= \Im \int_\T u_0 v_0 dx.
\]
Hence, there exists $N \in \N$ such that $\Im M (u_{0}^{(\l)}, v_0) \neq 0$ for $\l \ge N$.
Thus, we can take $\{(u_{0, n}, v_{0, n})\}_{n \in \N} = \{(u_{0}^{(n+N)}, v_0)\}_{n \in \N}$.

Next, we consider the case $(u_0, v_0) \notin \H^{s+\dl}(\T)$.
Without loss of generality, we assume $u_0 \notin H^{s+\dl}(\T)$.
We separate this case into the following three cases:
\begin{itemize}
\item
Case 1-1:
$P_+ u_0 \notin H^{s+\dl}(\T)$ and $P_- u_0 \notin H^{s+\dl}(\T)$;

\item
Case 1-2:
$P_+ u_0 \in H^{s+\dl}(\T)$ and $P_- u_0 \notin H^{s+\dl}(\T)$;

\item
Case 1-3:
$P_+ u_0 \notin H^{s+\dl}(\T)$ and $P_- u_0 \in H^{s+\dl}(\T)$.
\end{itemize}

Case 1-1:
We can take $\{(u_{0,n}, v_{0,n})\}_{n \in \N} = \{(u_{0}, v_0)\}_{n \in \N}$.

Case 1-2:
We set
\[
u_{0, +}^{(\l)} (x) = u_0(x) + \frac 1\l \sum_{k \in 2^\N} k^{-s-\dl} e^{i k x}
\]
for $\l \in \N$.
The same argument as in the case $(u_0, v_0) \in \H^{s+\dl}(\T)$  shows that
we can take
$\{(u_{0,n}, v_{0,n})\}_{n \in \N} = \{(u_{0,+}^{(n+N)}, v_0)\}_{n \in \N}$ for some large $N \in \N$.

Case 1-3:
This case is similarly handled as in Case 1-2.

Case 2: Assume that $\Im M (u_0, v_0) = 0$.
We claim that there exists $\{(u_{0}^{(\l)}, v_{0}^{(\l)} ) \}_{\l \in \N} \subset \H^s(\T)$ such that
\begin{equation}
\| (u_0^{(\l)}, v_0^{(\l)}) - (u_{0}, v_{0}) \|_{\H^s} = \frac 1\l 
\quad \text {and} \quad \Im M (u_{0}^{(\l)}, v_{0}^{(\l)}) \neq 0.
\label{c:app}
\end{equation}

If $P_0 v_0 \neq 0$, we set
\[
(u_0^{(\l)}, v_0^{(\l)}) =
\left\{
\begin{aligned}
&\Big( u_0 + \frac 1\l, v_0 \Big), && \text {if $\Im P_0 v_0 \neq 0$}, \\
&\Big( u_0 + \frac i\l, v_0 \Big), && \text {if $\Im P_0 v_0 = 0$}.
\end{aligned}
\right.
\] 
Then, a direct calculation shows that
\[
\| (u_0^{(\l)}, v_0^{(\l)}) - (u_0 , v_0) \|_{\H^s} 
= \frac 1\l.
\]
Moreover, we obtain
\begin{equation*}
\Im M (u_0^{(\l)}, v_0^{(\l)})
=
\Im \int_\T \Big(u_0 v_0 + \frac 1\l v_0 \Big) \; dx 
= 
\frac 1\l \Im P_0 v_0
\neq 0
%\label{appv0}
\end{equation*}
when $\Im P_0 v_0 \neq 0$.
Similarly, we obtain
\[
\Im M (u_0^{(\l)}, v_0^{(\l)})
=
\Im \int_\T \Big( u_0 v_0 + \frac i\l v_0 \Big) \; dx 
= 
\frac 1\l \Re P_0 v_0
\neq 0
\]
provided with $\Im P_0 v_0 = 0$.
Since $P_0 v_0 \neq 0$,
$\{(u_0^{(\l)}, v_0^{(\l)})\}_{\l \in \N}$ satisfies \eqref{c:app}.


If $P_0 u_0 \neq 0$, define
\[
(u_0^{(\l)}, v_0^{(\l)}) =
\left\{
\begin{aligned}
&\Big( u_0, v_0 + \frac 1\l \Big), && \text {if $\Im P_0 u_0 \neq 0$}, \\
&\Big( u_0, v_0 + \frac i\l \Big), && \text {if $\Im P_0 u_0 = 0$}.
\end{aligned}
\right.
\] 
The same argument as in the case of $\Im P_0 v_0 \neq 0$ implies that
$\{(u_0^{(\l)}, v_0^{(\l)})\}_{\l \in \N}$ satisfies \eqref{c:app}.

When $P_0 u_0 = 0$ and $P_0 v_0 = 0$ hold,
we take $\{(u_0^{(\l)}, v_0^{(\l)})\}_{\l \in \N}$ as
\[
(u_0^{(\l)}, v_0^{(\l)}) = \Big( u_0 + \frac 1{\sqrt \l}, v_0 + \frac i {\sqrt \l} \Big).
\]
Then, we have 
\[
\| (u_0^{(\l)}, v_0^{(\l)}) - (u_0, v_0) \|_{\H^s} = \frac 1\l.
\]
Moreover, 
\[
\begin{aligned}
\Im \int_\T  u_0^{(\l)} v_0^{(\l)} dx 
&=
\Im \int_\T u_0 v_0 dx + \frac 1{\sqrt \l}\Im P_0 v_0 + \frac 1{\sqrt \l} \Re P_0 u_0 
+ \frac 1{\l}  
\\
&= \frac 1{\l} \neq 0.
\end{aligned}
\]
Therefore, \eqref{c:app} holds.

The same argument as in Case 1 with \eqref{c:app} yields that
for any $\l \in \N$,
there exists $\{(u_0^{(\l, m)}, v_0^{(\l, m)}) \}_{m \in \N}$ such that
the following conditions hold:
\begin{enumerate}
%\item $(u_{0,n}, v_{0,n}) \in \H^s(\T) \setminus \H^{s+1}(\T)$
\item $(u_0^{(\l, m)}, v_0^{(\l, m)})$ converges to $(u_0^{(\l)}, v_0^{(\l)})$ 
in $\H^s(\T)$ as $m \to \infty$;
\item $\Im M (u_0^{(\l, m)}, v_0^{(\l, m)}) \neq 0$ for any $m \in \N$;
\item $(P_\pm u_0^{(\l, m)}, P_\pm v_0^{(\l, m)}) \notin \H^{s+\dl}(\T)$ for any $m \in \N$.
\end{enumerate}
In particular,
it follows from the proof of the Case 1 that
\begin{equation}
\| (u_0^{(\l, m)}, v_0^{(\l, m)}) - (u_0^{(\l)}, v_0^{(\l)}) \|_{\H^s}
\les
\frac 1m.
\label{app3}
\end{equation}
%We take $\{ (u_{0, n}, v_{0, n}) \}_{n \in \N}$ as 
%$ \{ (u_{0,n}^{(n)}, v_{0,n}^{(n)}) \}_{n \in \N}$.
The triangle inequality, \eqref{c:app}, and \eqref{app3} imply that
\[
\begin{aligned}
&\| (u_0^{(n,n)}, v_0^{(n,n)}) - (u_0, v_0) \|_{\H^s}
\\
&\le
\| (u_0^{(n, n)}, v_0^{(n, n)}) - (u_0^{(n)}, v_0^{(n)}) \|_{\H^s}
+ \| (u_0^{(n)}, v_0^{(n)}) - (u_0, v_0) \|_{\H^s}
\\
&\les \frac 1n.
\end{aligned}
\]
Thus, we can take $\{ (u_{0, n}, v_{0, n}) \}_{n \in \N}$ as 
$ \{ (u_0^{(n,n)}, v_0^{(n,n)}) \}_{n \in \N}$.
This concludes the proof.
\end{proof}

\end{document}
